\section{Bounds for the numerical enclosures}\label{or:threshold}

We prove the bounds used in Appendix~\ref{app:threshold}, using
its notation and the smoothing inequality~\eqref{eq:thr-fixed}.
These locate the threshold of Theorem~\ref{thm:threshold} and
certify the critical derivatives.
Write $\Phi_N(t)=\sum_{n\le N}c_n(it)^n$ for a partial sum of the
moment series.

\subsection{Bounds at a fixed memory}

\begin{lemma}[moment bounds]\label{lem:thr-moments}
Let $N\ge3$.
\begin{enumerate}[(a)]
\item Put $A=((N+1)a-1)/N$. If $m_j\le AB^j$ for $1\le j\le N$, then
$m_n\le AB^n$ for all $n\ge1$.
\item Let $S\subset\mathbb C$ be compact, with $\frac12<\alpha\le\operatorname{Re}z$,
$\lvert\operatorname{Im}z\rvert\le\rho$ and $\lvert z\rvert\le1$ for $z\in S$.
Define $m_n(z)$ by the recursion with $z$ in place of $a$, and
$c_n(z)=m_n(z)/\Gamma(1+nz)$. Then every $m_n$ is analytic near $S$. Put
$A=((N+1)\alpha-1)/N$. If $\lvert m_j\rvert\le AB^j$ on $S$ for $1\le j\le N$,
then for all $n\ge1$ and $z\in S$
\[
 \lvert m_n(z)\rvert\le AB^n,\qquad
 \lvert c_n(z)\rvert\le\frac{\tilde A\tilde B^n}{\Gamma(1+n\alpha)},\qquad
 \tilde A=Ae^{\rho^2/(2\alpha^2)},\quad \tilde B=Be^{\rho^2/(2\alpha)} .
\]
\end{enumerate}
\end{lemma}

\begin{proof}
For (a) we use induction on $n>N$. As $\theta_j\le1$ and
$na-\theta_n\ge na-1>0$, the bound for smaller indices gives
$m_n\le(n-1)A^2B^n/(na-1)$. The ratio $(na-1)/(n-1)$ increases in $n$, since
its derivative is $(1-a)/(n-1)^2$, and it equals $A$ at $n=N+1$. So
$m_n\le AB^n$.

For (b), $\operatorname{Re}(nz-\theta_n)\ge n\alpha-1>0$ on $S$, because
$\alpha\le1$. So $m_n$ is analytic near $S$, and the induction of (a) goes
through with $\lvert\theta_j\rvert\le1$. For $x>0$ and real $y$,
\[
 \Bigl\lvert\frac{\Gamma(x)}{\Gamma(x+iy)}\Bigr\rvert^2
 =\prod_{k\ge0}\Bigl(1+\frac{y^2}{(x+k)^2}\Bigr)
 \le\exp\Bigl(\frac{y^2}{x^2}+\frac{y^2}x\Bigr)
\]
by~\cite[5.8.3]{dlmf}. We take $x=1+n\operatorname{Re}z\ge1+n\alpha$ and
$\lvert y\rvert\le n\rho$. Then $y^2/x^2\le\rho^2/\alpha^2$ and
$y^2/x\le n\rho^2/\alpha$. Finally
$\Gamma(1+n\operatorname{Re}z)\ge\Gamma(1+n\alpha)$, as $\Gamma$ increases on
$[1.47,\infty)$.
\end{proof}

\begin{lemma}[truncation]\label{lem:thr-trunc}
Let $0<\alpha\le1$, and suppose that
$\lvert c_n\rvert\le\tilde A\tilde B^n/\Gamma(1+n\alpha)$ for $n>N$. Let $t>0$,
$0\le k\le N$ and $n_0=N+1$, and put
\[
 u=\frac{\tilde A\tilde B^{n_0}n_0^kt^{n_0-k}}{\Gamma(1+\alpha n_0)},\qquad
 r=\tilde Bt\Bigl(1+\frac1{n_0}\Bigr)^k
 \frac{\Gamma(1+\alpha n_0)}{\Gamma(1+\alpha n_0+\alpha)} .
\]
If $r<1$, then
$\bigl\lvert\frac{d^k}{dt^k}\sum_{n>N}c_n(it)^n\bigr\rvert\le\tau_k(t):=u/(1-r)$,
and $\tau_k$ increases on $(0,t]$.
\end{lemma}

\begin{proof}
The $k$th derivative of $(it)^n$ has modulus at most $n^kt^{n-k}$. The ratio
of consecutive terms of the majorant
$\tilde A\tilde B^nn^kt^{n-k}/\Gamma(1+\alpha n)$ is
$\tilde Bt(1+1/n)^k\Gamma(1+\alpha n)/\Gamma(1+\alpha n+\alpha)$. It decreases
in $n$, because $\log\Gamma$ is convex. So the majorant is at most a
geometric series with first term $u$ and ratio $r$. Both $u$ and $r$ increase
in $t$.
\end{proof}

\begin{lemma}[frequency cells]\label{lem:thr-cells}
Let $\nu$ be even, $t_c>0$ and $\lvert\delta\rvert\le h$. Then
\[
 \lvert\phi(t_c+\delta)\rvert\le\sum_{k<\nu}\lvert\phi^{(k)}(t_c)\rvert
 h^k/k!+c_\nu h^\nu.
\]
\end{lemma}

\begin{proof}
This is Taylor's formula, with
$\lvert\phi^{(\nu)}\rvert\le\E\lvert \mathcal L_a\rvert^\nu=\E \mathcal L_a^\nu=\nu!\,c_\nu$.
\end{proof}

\begin{lemma}[Fourier tail]\label{lem:thr-tail}
Let $0<T_2<T_3$ with $T_3\ge2^aT_2$, and let $\beta>3$ with $a\beta>1$.
Cover $[0,T_3]$ by cells $[l_j,r_j]$ with $l_0=0$ and with $T_2$ among the
end points, and let $s_j\ge\sup_{[l_j,r_j]}\lvert\phi\rvert$. Let $C$ be the
largest $s_jr_j^\beta$ over the cells that meet $[T_2,T_3]$, and put
\[
 R=T_2^{-1/a}\sum_{r_j\le T_2}s_j
 \Bigl(1-\frac12\Bigl(\frac{r_j}{T_2}\Bigr)^{1/a}\Bigr)^{-a\beta}
 \bigl(r_j^{1/a}-l_j^{1/a}\bigr)+\frac{2^{a\beta}CT_2^{-\beta}}{a\beta-1} .
\]
If $R<1$, then $\lvert\phi(t)\rvert\le Ct^{-\beta}$ for every $t\ge T_3$.
\end{lemma}

\begin{proof}
Let $b=s_j$ on the cells in $[0,T_3)$ and $b(t)=Ct^{-\beta}$ for $t\ge T_3$.
On a cell that meets $[T_2,T_3]$ we have $s_j\le Cr_j^{-\beta}\le Ct^{-\beta}$.
So $b(t)\le Ct^{-\beta}$ for $t\ge T_2$, and $\lvert\phi\rvert\le b$ on
$[0,T_3]$. Let $J(t)=\int_0^1b(tv^a)\,b(t(1-v)^a)\,dv$.

We show that $J(t)\le Rb(t)$ for $t\ge T_3$. By symmetry $J(t)$ is twice the
integral over $v\le\frac12$. There $t(1-v)^a\ge2^{-a}T_3\ge T_2$, so
$b(t(1-v)^a)\le b(t)(1-v)^{-a\beta}$. Hence $J(t)/b(t)$ is at most
$2\int_0^{1/2}b(tv^a)(1-v)^{-a\beta}\,dv$. Put $v_0=(T_2/t)^{1/a}\le\frac12$.
On $v<v_0$ we substitute $u=tv^a$, and this part becomes
\[
 \frac2at^{-1/a}\int_0^{T_2}s(u)\bigl(1-(u/t)^{1/a}\bigr)^{-a\beta}
 u^{1/a-1}\,du .
\]
Here $t^{-1/a}\le\frac12T_2^{-1/a}$ and $(u/t)^{1/a}\le\frac12(u/T_2)^{1/a}$.
The weight increases in $u$, so on each cell we take it at $r_j$. As
$\int_l^ru^{1/a-1}du=a(r^{1/a}-l^{1/a})$, this part is at most the first term
of $R$. On $v\ge v_0$ we use $b(tv^a)\le Ct^{-\beta}v^{-a\beta}$ and
$(1-v)^{-a\beta}\le2^{a\beta}$. This part is at most
$2^{1+a\beta}T_2^{-\beta}(T_2/t)^{1/a}C/(a\beta-1)$, which is at most the
second term of $R$.

Let $\varsigma$ be the supremum of the $\tau\ge T_3$ with
$\lvert\phi\rvert\le b$ on $[0,\tau]$, and suppose that $\varsigma<\infty$. By
continuity $\lvert\phi\rvert\le b$ on $[0,\varsigma]$. The arguments in
\eqref{eq:thr-fixed} at $t=\varsigma$ lie in $[0,\varsigma]$, so
$\lvert\phi(\varsigma)\rvert\le J(\varsigma)\le Rb(\varsigma)<b(\varsigma)$.
As $\lvert\phi\rvert$ and $b$ are continuous on $[T_3,\infty)$, we get
$\lvert\phi\rvert<b$ near $\varsigma$, against the choice of $\varsigma$.
\end{proof}

This is an explicit form of Step~4 in the proof of~\cite[Theorem~2.5]{glr2026},
with computed constants.

\begin{lemma}[the second derivative]\label{lem:thr-fpp}
If $\lvert\phi(t)\rvert\le Ct^{-\beta}$ for $t\ge T_3$, with $\beta>3$, then
$g_a$ has two continuous derivatives and
\[
 g_a''(0)=-\frac1\pi\Bigl(\sum_{\substack{n\le N\\ n\text{ even}}}
 (-1)^{n/2}c_n\frac{T_3^{n+3}}{n+3}+E\Bigr),\qquad
 \lvert E\rvert\le\frac{T_3^3}3\tau_0(T_3)+\frac{CT_3^{3-\beta}}{\beta-3} .
\]
\end{lemma}

\begin{proof}
As $t^2\operatorname{Re}\phi(t)$ is integrable, Fourier inversion gives
$g_a(x)=\frac1{2\pi}\int e^{-itx}\operatorname{Re}\phi(t)\,dt$, and dominated
convergence allows two derivatives under the
integral~\cite[Section~3.3]{durrett2019}. So
$g_a''(0)=-\frac1\pi\int_0^\infty t^2\operatorname{Re}\phi(t)\,dt$. On
$[0,T_3]$ we replace $\phi$ by $\Phi_N$, whose real part is
$\sum_{n\text{ even}}(-1)^{n/2}c_nt^n$, and integrate term by term. By
Lemma~\ref{lem:thr-trunc} the error is at most
$\int_0^{T_3}t^2\tau_0(t)\,dt\le\frac13T_3^3\tau_0(T_3)$. The integral over
$t>T_3$ is at most $CT_3^{3-\beta}/(\beta-3)$.
\end{proof}

\subsection{Bounds on memory intervals}

Near the threshold, the series for $\phi$ cancels terms of size
about $10^{60}$. Direct interval evaluation would need intervals
in $a$ narrower than about $10^{-65}$. We instead expand in $a$
about one exact point, so the cancellation occurs before we bound
the remainder.

\begin{lemma}[Cauchy remainder]\label{lem:thr-taylor}
Let $h$ be analytic near the closed disc $\lvert z-a_0\rvert\le\rho$, with
$\lvert h\rvert\le M$ there and Taylor coefficients $h_i$ at $a_0$. If $0<\vartheta<1$
and $\lvert a-a_0\rvert\le \vartheta\rho$, then
$\lvert h(a)-\sum_{i\le K}h_i(a-a_0)^i\rvert\le M\vartheta^{K+1}/(1-\vartheta)$.
\end{lemma}

\begin{proof}
Cauchy's estimate gives $\lvert h_i\rvert\le M\rho^{-i}$. Sum the geometric
tail.
\end{proof}

Let $I=[a_0-w,a_0+w]$, $0<\vartheta<1$ and $\rho=w/\vartheta$, and let $S$ be the square
$a_0+[-\rho,\rho]+i[-\rho,\rho]$, which contains the disc of
Lemma~\ref{lem:thr-taylor}. Suppose that $\alpha=a_0-\rho>\frac12$, that
$\lvert z\rvert\le1$ on $S$, and that $\tilde A,\tilde B$ are as in
Lemma~\ref{lem:thr-moments}(b). Let $\tau_k$ be the bound of
Lemma~\ref{lem:thr-trunc} with these constants. It holds for every real
$a\in I$. For a cell with center $t_c$ write
$\partial_t^k\Phi_N(z,t_c)=E_k(z)+iO_k(z)$, where $E_k$ is the sum over even
$n$ and $iO_k$ the sum over odd $n$. Both $E_k$ and $O_k$ are real for real
$z$. Let $E_{k,i}$ and $O_{k,i}$ be their Taylor coefficients at $a_0$, and
put $M_k=\sum_{n\le N}\sup_S\lvert c_n\rvert\,n(n-1)\cdots(n-k+1)\,t_c^{n-k}$.

\begin{proposition}[enclosure on an interval]\label{prop:thr-interval}
\begin{enumerate}[(a)]
\item For every $a\in I$ and every $t$ in the cell with center $t_c$ and
half-width $h$,
\[
 \lvert\phi_a(t)\rvert\le\sum_{k<12}\Bigl(\sum_{i\le K}w^i
 \sqrt{E_{k,i}^2+O_{k,i}^2}+\frac{M_k\vartheta^{K+1}}{1-\vartheta}+\tau_k(t_c)\Bigr)
 \frac{h^k}{k!}+h^{12}\sup_S\lvert c_{12}\rvert .
\]
\item Let $T_2$, $T_3\ge2^{a_0+w}T_2$, $\beta$ and the cells be as in
Lemma~\ref{lem:thr-tail}, with $s_j$ the bound of (a), and suppose that $R<1$
for every $a\in I$. Put
\[
 F(z)=\sum_{n\le N,\,n\text{ even}}(-1)^{n/2}c_n(z)T_3^{n+3}/(n+3).
\]
Let $F_i$ be its Taylor coefficients at $a_0$, and let $M_F$ be
the same sum with $\sup_S\lvert c_n\rvert$ in place of $(-1)^{n/2}c_n(z)$. Then
for every $a\in I$
\[
 \Bigl\lvert g_a''(0)+\frac1\pi\sum_{i\le K}F_i(a-a_0)^i\Bigr\rvert\le\eta
 =\frac1\pi\Bigl(\frac{M_F\vartheta^{K+1}}{1-\vartheta}+\frac{T_3^3}3\tau_0(T_3)
 +\frac{CT_3^{3-\beta}}{\beta-3}\Bigr),
\]
and $a\mapsto g_a''(0)$ is continuous on $I$.
\end{enumerate}
\end{proposition}

\begin{proof}
For (a) we apply Lemma~\ref{lem:thr-cells} with $\nu=12$, and
Lemma~\ref{lem:thr-trunc} to replace $\phi_a^{(k)}$ by $\partial_t^k\Phi_N$.
For real $a$ we have
$\lvert\partial_t^k\Phi_N(a,t_c)\rvert=\lvert(E_k(a),O_k(a))\rvert$.
Lemma~\ref{lem:thr-taylor} for $E_k$ and $O_k$, with
$\sup\lvert E_k\rvert+\sup\lvert O_k\rvert\le M_k$ on $S$, and the triangle
inequality in $\mathbb R^2$ give the bracket. Also
$c_{12}(a)\le\sup_S\lvert c_{12}\rvert$.

For (b), $T_3\ge2^aT_2$ for every $a\in I$, so Lemma~\ref{lem:thr-tail} applies
at each $a$. Then Lemma~\ref{lem:thr-fpp} and Lemma~\ref{lem:thr-taylor} for
$F$ give the enclosure. By Lemma~\ref{lem:thr-moments}(b) the series for
$\phi_a(t)$ converges uniformly for $a\in I$ and $t$ in compact sets. So
$\phi_a(t)$ is continuous in $a$. Moreover $t^2\lvert\phi_a(t)\rvert$ is at
most $t^2$ on $[0,T_3]$ and $Ct^{2-\beta}$ beyond, uniformly in $a\in I$.
Dominated convergence gives the continuity.
\end{proof}

\subsection{Critical derivative bounds}
\label{or:transition-unique}

To compute the memory derivative and critical constants, we
differentiate the smoothing equation and then bound the tail. The error bound in Lemma~\ref{lem:thr-fpp}
alone gives no bound on the derivative of that error.

\begin{lemma}[differentiated smoothing bound]\label{lem:thr-parameter-tail}
Fix a compact memory interval $I\subset(1/2,1)$. Suppose that the
cell certificate of Lemma~\ref{lem:thr-tail} gives, uniformly on $I$,
$f(t):=|\phi_a(t)|\le Ct^{-\beta}$ for $t\ge T_2$, and
\[
 H_\beta(t):=2\int_0^{1/2}f(tv^a)(1-v)^{-a\beta}\,dv\le R<1
 \quad(t\ge T),\qquad T\ge2^aT_2.
\]
Let $\beta>4$, $\gamma=\beta-1$, $a\gamma>1$, and
$u=|\partial_t\phi_a|$, $w=|\partial_a\phi_a|$.
Suppose that $u\le M$ everywhere, $w\le M_a$ on $[0,T_2]$, and
cells give $u(t)<D_1t^{-\beta}$ and $w(t)<D_at^{-\gamma}$ on $[T_2,T]$.
Put
\begin{align*}
 A_1&=\frac{MC2^{a(\beta-1)}}{a+1},&
 J&=2C2^{a\beta}T^{\gamma-\beta-1/a},\\
 \kappa&=R+\frac{JT_2^{1/a-\gamma}}{a\gamma-1},&
 A_0&=JM_aT_2^{1/a},\\
 A_S&=MC2^{a(\beta-1)}
 \left\{\frac{\log2}{a+1}+\frac1{(a+1)^2}\right\}
       +D_1(\log2)R+CR/T.
\end{align*}
If $A_1+RD_1<D_1$ and $A_0+A_S+\kappa D_a<D_a$, with
$\kappa<1$, then the two derivative bounds hold for every $t\ge T_2$.
Consequently, in the interior of $I$,
\begin{equation}\label{eq:thr-memory-derivative}
 \frac{d}{da}g_a''(0)=-\frac1\pi\int_0^\infty
 t^2\operatorname{Re}\partial_a\phi_a(t)\,dt,
 \qquad
 \left|\frac1\pi\int_T^\infty t^2
 \operatorname{Re}\partial_a\phi_a(t)\,dt\right|
 \le\frac{D_aT^{3-\gamma}}{\pi(\gamma-3)}.
\end{equation}
\end{lemma}

\begin{proof}
The complex majorant in Lemma~\ref{lem:thr-moments}(b) makes the
moment series locally uniformly convergent in the frequency and memory.
Cauchy's estimate on a smaller parameter disc gives the same conclusion
for its memory derivative and every fixed frequency derivative.
We may therefore differentiate the series for $\phi_a$ term by term.
The smoothing equation, with $p=(1+a)/2$, is
\[
 \phi_a(t)=\int_0^1\phi_a(tv^a)
 \{p\phi_a(t(1-v)^a)+(1-p)\phi_a(-t(1-v)^a)\}\,dv.
\]
On compact frequency sets, its differentiation under the integral is
justified by local boundedness of these derivatives and boundedness of
$v^{\inf I}|\log v|$. Conjugate symmetry gives even moduli.
Writing $x=tv^a$ and $y=t(1-v)^a$, we obtain
\begin{align*}
 u(t)&\le\int_0^1\{v^au(x)f(y)+(1-v)^af(x)u(y)\}\,dv,\\
 w(t)&\le\int_0^1\{w(x)f(y)+f(x)w(y)\}\,dv+S(t),\\
 S(t)&\le\int_0^1\{x|\log v|u(x)f(y)
             +y|\log(1-v)|f(x)u(y)+f(x)f(y)\}\,dv.
\end{align*}
The final term includes the derivative of $p$. Its coefficient is one because $p'=1/2$ and $|\phi_a(y)-\phi_a(-y)|\le2f(y)$.

Split each integral at $1/2$ and reflect its second half. For the
frequency derivative, the smaller-argument term is at most
$A_1t^{-\beta}$, by $u\le M$ and
\[
 2\int_0^{1/2}v^a(1-v)^{-a\beta}\,dv
 \le\frac{2^{a(\beta-1)}}{a+1}.
\]
The larger-argument term is at most $RD_1t^{-\beta}$.
For the memory derivative, split the smaller-argument term again at
$v_0=(T_2/t)^{1/a}$. Its bound is
\[
 2C2^{a\beta}t^{-\beta-1/a}
 \left\{M_aT_2^{1/a}+
 \frac{D_aT_2^{1/a-\gamma}}{a\gamma-1}\right\}.
\]
The larger-argument term is at most $RD_at^{-\gamma}$, since
$H_\gamma\le H_\beta$. As $\gamma-\beta-1/a<0$ and $t\ge T$,
these two terms are bounded by $(A_0+\kappa D_a)t^{-\gamma}$.
In $S$, the smaller logarithmic term uses
\[
 \int_0^{1/2}v^a(-\log v)\,dv
 =2^{-a-1}\left\{\frac{\log2}{a+1}+\frac1{(a+1)^2}\right\};
\]
the larger uses $-\log(1-v)\le\log2$ and $H_\beta\le R$.
The mixing-weight term is at most $CRt^{-\beta}$.
Thus $S(t)\le A_St^{-\gamma}$.

To extend the bounds beyond the cells, maximize the continuous weighted derivative on any $[T_2,U]$.
If its maximum lies in $[T_2,T]$, use the cell bound. Otherwise repeat
the preceding inequalities with that maximum in place of the proposed
constant, and rearrange the strict contraction inequality. This first
proves the frequency bound, then the memory bound, independently of $U$.
The integrable uniform bound $t^2w(t)\le D_at^{2-\gamma}$ now permits
differentiation of the Fourier formula for $g_a''(0)$ on the interior
of $I$ and gives~\eqref{eq:thr-memory-derivative}.
\end{proof}

The computation uses
\[
 I=[0.66881285,0.66881287],\quad T_2=6,\quad T=10,
 \quad\beta=8,\quad\gamma=7,
\]
and the conservative constants
\[
 C=17/1000,\quad R=83/1000,\quad M=2,
 \quad D_1=2,\quad M_a=D_a=4.
\]
It checks $A_1<1$, $\kappa<1/10$ and $A_0+A_S<2$, so both
contractions have strict slack. The resulting tail bound is
$1/(10000\pi)<1/30000$.
The certified zero bracket lies strictly inside $I$.

For the finite-frequency integral, let $w_0$ be the halfwidth of $I$,
$\rho>w_0$ the Cauchy radius and $\vartheta=w_0/\rho$.
Differentiating the finite Taylor polynomial in memory leaves the
coefficient remainder factor
\[
 \frac1\rho\sum_{i>K}i \vartheta^{i-1}
 =\frac{\vartheta^K\{(K+1)-K\vartheta\}}{\rho(1-\vartheta)^2}.
\]
The omitted moment series has derivative bound
$\tau_k(t)/(\rho-w_0)$ by Cauchy's estimate.
On each frequency cell we use a Taylor expansion of order $48$, with its remainder bounded by the complex coefficient majorant. (At complex memory we cannot use a probability moment.)
The signed interval checks give
\begin{equation}\label{eq:thr-final-derivatives}
 0.68<\frac{d}{da}g_a''(0)<0.70,\qquad
 -0.09<g_a^{(4)}(0)<-0.07
 \quad(a\in\operatorname{int}I).
\end{equation}
The second enclosure uses the same Fourier bound with weight $t^4$, whose tail is $CT^{-3}/(3\pi)$.
The file \texttt{certify\_transition.py} performs these checks, including
the initial derivative-cell bounds. The recorded run uses $1100$ moments
and $1600$-bit precision under \texttt{python -O}.
Each check uses \texttt{require}, so it also runs with optimization enabled.

At the unique zero from Theorem~\ref{thm:threshold}, evenness
and the second inequality in~\eqref{eq:thr-final-derivatives} give
\[
 g_{a_*}'(x)=\tfrac16g_{a_*}^{(4)}(0)x^3+o(x^3)<0
 \quad\text{for all sufficiently small }x>0.
\]
The fourth derivative is nonzero, so the critical top is quartic.
The analytic theorem already gives the unimodal critical shape
and excludes other curvature zeros on $(1/2,1)$.

\paragraph{The critical constants.}
For Table~\ref{tab:critical-constants} we apply the same estimates on
$I_*=[z_--10^{-14},z_-+10^{-14}]$, where
$z_-=1838418790980431/2748779069440000$ is the lower endpoint of the
certified threshold bracket. The entire bracket lies inside $I_*$.
Take
\[
 \begin{gathered}
 T_2=10,\quad T=16,\quad\beta=12,\quad\gamma=11,\\
 C=1/1000,\quad R=1/20,\quad M=D_1=2,\quad M_a=D_a=4.
 \end{gathered}
\]
The base Fourier bound, initial derivative cells and both strict
contraction inequalities hold uniformly on $I_*$. We use $2000$ moments,
$2600$-bit ball arithmetic, Cauchy ratio $1/10000$, frequency halfwidth
$1/32$ and frequency Taylor order $48$.
The omitted tails for $C_*$, $B_*$ and $H_*$ are bounded respectively by
\[
 \frac{D_aT^{3-\gamma}}{\pi(\gamma-3)},\qquad
 \frac{CT^{5-\beta}}{\pi(\beta-5)},\qquad
 \frac{CT^{2-\beta}}{\pi(\beta-2)}.
\]
The finite-frequency integrals use the same parameter and moment
remainders as above, with weights $t^2$, $t^4$ and $t$.
The bounds hold throughout the interval, so they apply at $a_*$
without needing its exact value. Outward rounding then gives the
two derived coefficients in the table. The record and rerun command are
in Appendix~\ref{app:data}.

